\documentclass{subfiles}
\begin{document}\label{sec:DC3}
    Before we describe the algorithm let us make a remark: suffix arrays are 
        not the only option available when it comes to sorting suffixes.
        One could also use a \gls{st}, whose construction has been proposed 
        in several papers; see \cite{farach1997,farachM1996,mccreight1976,ukkonen1995} 
        for instance. The only issue is that, ST are space inefficient.

    The algorithm we describe here was proposed by K{\"a}rkk{\"a}inen et al. in 
        \cite{karkkainenSB2006}. Broadly speaking, the problem of sorting 
        suffixes of a string of length \(n\), 
        is recursively reduced to that of sorting the suffixes of a string of
        length \(\frac{2}{3}n\).

    The algorithm makes the following trivial assumption:
        the string is defined over a integer alphabet.
        Using the above assumption, we use \([i, j]\) as a shorthand for \(\set{1, \ldots, j}\).

    Let \(T[0, n]\) be the text to process. Apppend a end of string symbol \((\textdollar)\)
        to it. For \(k = 0, 1, 2\), define
        \[
            B_{k} = \set{i \in [0, n + 1]}[i \equiv k \mod{3}].
        \]

    Let \(R_{k}\) denote the string \(T[k, n]\), for \(k = 1, 2\). 
        Consider \(R_{k}\)'s characters as triplets of the form \([t_{i}, t_{i + 1}, t_{i + 2}]\),
        eventually padded by and end of string symbol, for \(i \in B_{k}\).
        Let \(R = R_{1}R_{2}\). It can be shown that there exist a correspondence 
        between the suffixes of \(R\) an the suffixes of \(S\).
    
    To sort \(R\)'s suffixes, radix sort its characters (the triplets) and replace 
        each with its rank; denote this new string with \(R'\).
        If all characters are distinct, the ranks gives the order of the suffixes;
        otherwise sort \(R'\) suffixes using the DC3 algorithm.

    After the suffixes have been computed, for \(i \in B_{1} \cup B_{2}\), let 
        \(rank(S_{i})\) denote the rank of suffix \(S_{i}\). For \(i \in B_{0}\),
        \(rank(S_{i})\) is undefined.

    Note that for \(i \in B_{0}, S_{i + 1}\) has been computed.
        Thus, if we represent each \(S_{i}\) as the pair \((t_{i}, rank(S_{i + 1}))\),
        it follows that for all \(i, j \in B_{0}\),
        \[
            S_{i} \le S_{j} \iff (t_{i}, rank(S_{i + 1}) \le (t_{j}, rank(S_{j + 1})). 
        \]
    
    At this point all suffixes whose starting position belongs to \(B_{1} \cup B_{2}\)
        have been sorted, and we have a way to sort suffixes starting at a position 
        in \(B_{0}\). What we are left with is merging the two. For the sake of simplicity,
        let \(j \in B_{0}\) and \(i \in B_{1} \cup B_{2}\). Then, to compare \(S_{i}\)
        and \(S_{j}\) we do the following:
        \begin{itemize}
            \item if \(j \in B_{1}\), compare \(S_{i}\) and \(S_{j}\) as both \(i, j \in B_{0}\).
                That is
                \[
                    S_{i} \le S_{j} \iff (t_{i}, rank(S_{i + 1})) \le (t_{j}, rank(S_{j + 1})).
                \]
            \item if \(j \in B_{2}\), compare \(t_{i}\) and \(t_{j}\), if these are equal 
                compare \(t_{i + 1}\) and \(t_{j + 1}\) and if this are equal compare 
                \(rank(S_{i + 2})\) and \(rank(S_{j + 2})\). That is, 
                \[
                    S_{i} \le S_{j} \iff
                        (t_{i}, t_{i + 1}, rank(S_{i + 2})) \le (t_{j}, t_{j + 1}, rank(S_{j + 2})).
                \]
        \end{itemize}

    Computationally speaking, apart from the recursive call, 
        all steps take linear time. The recursive call is on a string of length 
        \(\frac{2}{3}n\). Therefore the overall time complexity is given by the recurrence
        \(T(n) = T\left\lparen\frac{2}{3}n\right\rparen + \bigO{n}\), whose solution by the Master theorem is 
        \(\bigO{n}\) time.

    \begin{example*}%             0 1 2 3 4 5 6 7 8 9 10 11 12 13
        Consider the string \(w = m a t h i s a w e s  o  m  e  \$\). We get the buckets:
        \[
            B_{0} = \set{0, 3, 6, 9, 12},
            B_{1} = \set{1, 4, 7, 10, 13}, 
            B_{2} = \set{2, 5, 8, 11}
        \]

        For \(k = 1, 2\), we get 
            \[\begin{aligned}
                R_{1} & = [a,t,h][i,s,a][w,e,s][o,m,e][\$,\$,\$] \\
                R_{2} & = [t,h,i][s,a,w][e,s,o][m,e,\$]
            \end{aligned}\]
            thus, 
            \[\begin{aligned}
                R = & R_{1}R_{2} \\
                  = & [a,t,h][i,s,a][w,e,s][o,m,e][\$,\$,\$][t,h,i][s,a,w][e,s,o][m,e,\$].
                       %S_{1}  S_{4}  S_{7}  S_{10}  S_{13}   S_{2}  S_{5}  S_{8}  s_{11}
            \end{aligned}\]

        Using the radix sort on \(R\)'s characters we get
            \[
               R' = (1, 3, 8, 5, 0, 7, 6, 2, 4).
            \]
            since no two are equal, the suffixes are sorted. 
            What we are left is computing the suffixes for the bucket \(B_{0}\).
            By using the notation above, these are 
            \[\begin{aligned}
                S_{0} & = (m, rank(S_{1})) = (m, 1) \\ 
                S_{3} & = (h, rank(S_{4})) = (h, 3) \\ 
                S_{6} & = (a, rank(S_{7})) = (a, 8) \\ 
                S_{9} & = (s, rank(S_{10})) = (s, 5) \\ 
                S_{12} & = (e, rank(S_{13})) = (e, 0).
            \end{aligned}\]
            Lastly, comparing them with the previously computed suffixes, 
            the suffix array for \(w\) is 
            \[
                SA_{w} = (13, 1, 6, 12, 8, 3, 4, 0, 11, 10, 5, 9, 2, 7)
            \]
            since \([\$,\$,\$]\) represents the end of string symbol, 
            we can remove \(S_{13}\). 
    \end{example*}
\end{document}
